\section{实践示例与指令模板}

\subsection{Prompt Template 演练}
一个理想的 inference-time prompt 包含 context、task difficulty、example CoT、instrumentation hints 以及 final instruction。例如下面的模板在多步骤 reasoning 中较常使用：

\begin{lstlisting}[language=Python, caption={Inference-time prompt skeleton}]
System: You are a high-precision reasoning assistant.

Input context: {problem description}
Difficulty: {mark as "hard"/"critical"}
Examples:
  - Example 1 CoT: ...
  - Example 2 CoT: ...
Instruction: Think step-by-step. If uncertain, enumerate two next steps and sample multiple candidates.
\end{lstlisting}

\begin{importantbox}{提示模板的落地技巧}
把 difficulty 作为 metadata 绑定在 prompt 里，便于 sampling policy 自动调整 temperature 与 token limit；用例示例表达 “wrong → right” 对比，提醒模型要把 reasoning 拓展到深度。
\end{importantbox}

\subsection{Self-Consistency Orchestration}
在 sampling 阶段，先用 high-temperature 不断生成 20~40 个 responses；在 search 阶段，把这些 responses 看作 seed，开启 beam/tree search 进一步延展 promising branches；在 iteration 阶段，运行 trace feedback 和 verifier 评分，再基于 majority vote/consistency 选出最终 answer。

\begin{knowledgebox}{Orchestration checklist}
1. Sampling：记录 candidate count + diversity；2. Search：track branch depth, stuck rate；3. Iteration：log trace iteration count, verification score, self-correction delta。
\end{knowledgebox}

\subsection{本章小结}
实践示例帮助把 abstract 的 inference-time 理念具体化：我们可以借助 prompt skeleton、sampling/search/iteration 三阶段 orchestrations，以及 instrumentation checklist 来把 pipeline 写成可复用的 template。

